\documentclass[12pt,a4paper]{article}
\usepackage{amsmath,amssymb,mathrsfs,tikz,times,pifont}
\usepackage{enumitem}
\newcommand\circitem[1]{%
\tikz[baseline=(char.base)]{
\node[circle,draw=gray, fill=red!55,
minimum size=1.2em,inner sep=0] (char) {#1};}}
\newcommand\boxitem[1]{%
\tikz[baseline=(char.base)]{
\node[fill=cyan,
minimum size=1.2em,inner sep=0] (char) {#1};}}
\setlist[enumerate,1]{label=\protect\circitem{\arabic*}}
\setlist[enumerate,2]{label=\protect\boxitem{\alph*}}
%%%::::::by chnini ameur :::::::%%%
\everymath{\displaystyle}
\usepackage[left=1cm,right=1cm,top=1cm,bottom=1.7cm]{geometry}
\usepackage{array,multirow}
\usepackage[most]{tcolorbox}
\usepackage{varwidth}
\tcbuselibrary{skins,hooks}
\usetikzlibrary{patterns}
%%%::::::by chnini ameur :::::::%%%
\newtcolorbox{exa}[2][]{enhanced,breakable,before skip=2mm,after skip=5mm,
colback=yellow!20!white,colframe=black!20!blue,boxrule=0.5mm,
attach boxed title to top left ={xshift=0.6cm,yshift*=1mm-\tcboxedtitleheight},
fonttitle=\bfseries,
title={#2},#1,
% varwidth boxed title*=-3cm,
boxed title style={frame code={
\path[fill=tcbcolback!30!black]
([yshift=-1mm,xshift=-1mm]frame.north west)
arc[start angle=0,end angle=180,radius=1mm]
([yshift=-1mm,xshift=1mm]frame.north east)
arc[start angle=180,end angle=0,radius=1mm];
\path[left color=tcbcolback!60!black,right color = tcbcolback!60!black,
middle color = tcbcolback!80!black]
([xshift=-2mm]frame.north west) -- ([xshift=2mm]frame.north east)
[rounded corners=1mm]-- ([xshift=1mm,yshift=-1mm]frame.north east)
-- (frame.south east) -- (frame.south west)
-- ([xshift=-1mm,yshift=-1mm]frame.north west)
[sharp corners]-- cycle;
},interior engine=empty,
},interior style={top color=yellow!5}}
%%%%%%%%%%%%%%%%%%%%%%%

\usepackage{fancyhdr}
\usepackage{eso-pic}         % Pour ajouter des éléments en arrière-plan
% Commande pour ajouter du texte en arrière-plan
\AddToShipoutPicture{
    \AtTextCenter{%
        \makebox[0pt]{\rotatebox{80}{\textcolor[gray]{0.7}{\fontsize{5cm}{5cm}\selectfont PGB}}}
    }
}
\usepackage{lastpage}
\fancyhf{}
\pagestyle{fancy}
\renewcommand{\footrulewidth}{1pt}
\renewcommand{\headrulewidth}{0pt}
\renewcommand{\footruleskip}{10pt}
\fancyfoot[R]{
\color{blue}\ding{45}\ \textbf{2025}
}
\fancyfoot[L]{
\color{blue}\ding{45}\ \textbf{Prof:M. BA}
}
\cfoot{\bf
\thepage /
\pageref{LastPage}}
\begin{document}
\renewcommand{\arraystretch}{1.5}
\renewcommand{\arrayrulewidth}{1.2pt}
\begin{tikzpicture}[overlay,remember picture]
\node[draw=blue,line width=1.2pt,fill=purple,text=blue,inner sep=3mm,rounded corners,pattern=dots]at ([yshift=-2.5cm]current page.north) {\begingroup\setlength{\fboxsep}{0pt}\colorbox{white}{\begin{tabular}{|*1{>{\centering \arraybackslash}p{0.28\textwidth}} |*2{>{\centering \arraybackslash}p{0.2\textwidth}|} *1{>{\centering \arraybackslash}p{0.19\textwidth}|} }
\hline
\multicolumn{3}{|c|}{$\diamond$$\diamond$$\diamond$\ \textbf{Lycée de Dindéfélo}\ $\diamond$$\diamond$$\diamond$ }& \textbf{A.S. : 2025/2026} \\ \hline
\textbf{Matière: Mathématiques}& \textbf{Niveau : 2nd}\textbf{S} &\textbf{Date: 20/12/2025} & \textbf{Durée : 4 heures} \\ \hline
\multicolumn{4}{|c|}{\parbox[c]{10cm}{\begin{center}
\textbf{{\Large\sffamily Composition Du 1$ ^\text{\bf ère} $ Semestre}}
\end{center}}} \\ \hline
\end{tabular}}\endgroup};
\end{tikzpicture}
\vspace{3cm}

\section*{\underline{Exercice 1 :} (5 points) }

\begin{enumerate}
    \item Rappeler la forme canonique de $P(x)=ax^2+bx+c$ \hfill \textbf{(0,5 pt)}
    
    \item Soient $P(x)= 2x^2-x+2$ et $Q(x)= 3x^2-2\sqrt{3}x+1$.
    \begin{enumerate}
        \item Donner la forme canonique de $P(x)$ et $Q(x)$ \hfill \textbf{(0,75 pt)}
        \item Factoriser si possible $P(x)$ et $Q(x)$ \hfill \textbf{(0,75 pt)}
    \end{enumerate}

    \item Résoudre dans $\mathbb{R}$ les équations et les inéquations suivantes :
    \begin{enumerate}
        \item $x^2 + 2x - 15 = 0$ \hfill \textbf{(1 pt)}
        \item $2x^2 - 9x + 4 \geq 0$ \hfill \textbf{(1 pt)}
        \item $6x^4 - x^2 - 5 \leq 0$ \hfill \textbf{(1 pt)}
    \end{enumerate}
\end{enumerate}

\section*{\underline{Exercice 2 :} ( 5 points) }

Résoudre dans $\mathbb{R}$ ou $\mathbb{R}^{2}$ les équations, inéquations et systèmes suivants :

\begin{enumerate}
    \item $2|x - 2|^2 + 7|x - 2| + 3 = 0$ \hfill \textbf{(1 pt)}
    
    \item $\dfrac{2x^2 + 3x - 5}{-x^2 + 3x + 10} \leq 0$ \hfill \textbf{(1 pt)}
    
    \item $\left(x + \dfrac{1}{x}\right)^2 = 4\left(x + \dfrac{1}{x}\right) - 3$ \hfill \textbf{(1 pt)}
    
    \item $\begin{cases} x + y = 3 \\ x^2 + 3xy + y^2 = 11 \end{cases}$ \hfill \textbf{(1 pt)}
    
    \item En utilisant la méthode de \textbf{Cramer}, résoudre le système : 
    $\begin{cases} 3x + 4y = 5 \\ 2x - 3y = 1 \end{cases}$ \hfill \textbf{(1 pt)}
\end{enumerate}


\newpage
\section*{\underline{Exercice 3 :} (6 points) }

Soit $ABC$ un triangle quelconque, $I$ le milieu de $[AC]$, $D$ le symétrique de $B$ par rapport à $C$.

\begin{enumerate}
    \item Faire une figure que l'on complétera au fur et à mesure de l'exercice. \textbf{(0,5 pt)}
    \item Justifier que $I = \text{bary}\{(A, 2); (C, 2)\}$ et $D = \text{bary}\{(B, -1); (C, 2)\}$. \textbf{(0,75 pt)}
    \item Soit $G = \text{bary}\{(A, 2); (B, -1); (C, 2)\}$.\\ Montrer que $G$ est le point d'intersection de $(AD)$ et $(BI)$. \textbf{(0,75 pt)}
    \item La droite $(CG)$ coupe $(AB)$ en $K$. Montrer que $A$ est le milieu de $[BK]$. \textbf{(0,75 pt)}
    \item Démontrer que $(DK)$ et $(AC)$ sont parallèles. \textbf{(0,5 pt)}
    \item Soit $E' = \text{bary}\{(A, 2); (B, 1)\}$.
    \begin{enumerate}
        \item Exprimer $\overrightarrow{E'I}$ et $\overrightarrow{ID}$ en fonction de $\overrightarrow{AB}$ et $\overrightarrow{AC}$. \textbf{(0,75 pt)}
        \item En déduire $E'$, $I$ et $D$ sont alignés. \textbf{(1 pt)}
    \end{enumerate}
    \item Que représente $G$ pour le triangle $BKD$ ? \textbf{(0,5 pt)}
    \item Soit $J$ le point d'intersection de $(BG)$ et $(KD)$. Que peut-on dire de $(AJ)$ et $(BD)$ ? Justifier votre réponse. \textbf{(0,5 pt)}
\end{enumerate}
\section*{\underline{Exercice 4 :} (4 points) }

\begin{enumerate}
\item Déterminer la mesure principale $\alpha \in ]-\pi ; \pi]$ des angles orientés suivants :

\begin{enumerate}
    \item $\alpha_1 = \dfrac{2027\pi}{4}$ \hfill \textbf{(0,5 pt)}
    \item $\alpha_2 = -\dfrac{31\pi}{6}$ \hfill \textbf{(0,5 pt)}
    \item $\alpha_3 = \dfrac{2026\pi}{7}$ \hfill \textbf{(0,5 pt)}
    \item $\alpha_4 = 40\pi$ \hfill \textbf{(0,5 pt)}
\end{enumerate}

\item Calculer les valeurs exactes suivantes : 
\begin{enumerate}
\item $\cos \left( \dfrac{97\pi}{4} \right)$ \hfill \textbf{(1 pts)}
\item $\sin \left( \dfrac{31\pi}{4} \right)$ \hfill \textbf{(1 pts)}
\end{enumerate}
\end{enumerate}

\end{document}
